%% =====================================================================
%%  config.tex — ไฟล์ตั้งค่าวิทยานิพนธ์ (กรอกข้อมูลที่นี่ที่เดียว)
%%  แก้ไขข้อมูลในไฟล์นี้ แล้ว compile main_th.tex หรือ main_en.tex
%% =====================================================================


%% ===================================================================
%%  Template Configuration)
%% ===================================================================

\def\MSUfontThai{16}      %% ขนาดฟอนต์ภาษาไทย — TH Sarabun New (pt)
\def\MSUfontEnglish{12}   %% ขนาดฟอนต์ภาษาอังกฤษ — Times New Roman (pt)

\headinglist
{ตาราง}
{Table}
{ภาพประกอบ}
{Figure}


%% ===================================================================
%%  1. ชื่อเรื่อง (Thesis Title)
%% ===================================================================
\thesistitles
{ชื่อวิทยานิพนธ์ภาษาไทย}            % Thai Title of the Thesis 
{English Title of the Thesis}   % English Title of the Thesis

%% ===================================================================
%%  2. ข้อมูลผู้วิจัย (Author Information)
%% ===================================================================
\authorprefix
{นาย}
{Mr.}

\thesisauthor
{ฉัตรมงคล อารีญาติ}
{Chatmongkol Areeyat}
{68010361011}

\advisor
{ศ. ดร. วรวัฒน์ เสงี่ยมวิบูล}
{ศาสตราจารย์ ดร. วรวัฒน์ เสงี่ยมวิบูล}
{Prof. Worawat Sa-ngiamwibool, Ph.D}
{Professor Worawat Sa-ngiamwibool, Ph.D}
{อาจารย์ที่ปรึกษาวิทยานิพนธ์หลัก}
{Advisor}

\coadvisorlist
{รศ. ดร. สุพรรณิกา วัฒนะ}
{รองศาสตราจารย์ ดร.สุพรรณิกา วัฒนะ}
{Assoc. Prof. Supannika Wattana, Ph.D}
{Associate Professor Supannika Wattana, Ph.D}
{ผศ. ดร. บัญชา วัฒนะ}
{ผู้ช่วยศาสตราจารย์ ดร.บัญชา วัฒนะ}
{Asst. Prof. Buncha Wattana, Ph.D}
{Assistant Professor Buncha Wattana, Ph.D}

\coadvisorName
{อาจารย์ที่ปรึกษาวิทยานิพนธ์ร่วม}
{Co-Advisor}

\committeeChair
{ผศ. ดร. ชัยยงค์ เสริมผล}
{ผู้ช่วยศาสตราจารย์ ดร.ชัยยงค์ เสริมผล}
{Asst. Prof. Chaiyong Soemphol, Ph.D}
{Assistant Professor Chaiyong Soemphol, Ph.D}
{ประธานกรรมการ}
{Chairman}

\committeelist  
% 1st commitee 
{ผศ. ดร. ณัฐวุฒิ สุวรรณทา}
{ผู้ช่วยศาสตราจารย์ ดร.ณัฐวุฒิ สุวรรณทา}
{Asst. Prof. Nattawoot Suwannata, Ph.D}
{Assistant Professor Nattawoot Suwannata, Ph.D}
%
% 2nd commitee 
{}
{}
{}
{}
%{รศ. ดร. ชลธี โพธิ์ทอง}
%{รองศาสตราจารย์ ดร.ชลธี โพธิ์ทอง}
%{Assoc. Prof. Chonlatee  Photong, Ph.D}
%{Associate Professor Chonlatee  Photong, Ph.D}

\committeeName
{กรรมการ}
{Committee}

\examiner
{ผศ. ดร. จักรพนธ์ อบมา}
{ผู้ช่วยศาสตราจารย์ ดร.จักรพนธ์ อบมา}
{Asst. Prof. Jagraphon Obma, Ph.D}
{Assistant Professor Jagraphon Obma, Ph.D}
{กรรมการผู้ทรงคุณวุฒิภายนอก}
{Examiner}

\deanEng
{รศ. ดร. จักรมาส เลาหวณิช}
{รองศาสตราจารย์ ดร.จักรมาส เลาหวณิช}
{Assoc. Prof. Juckamas Laohavanich, Ph.D}
{Associate Professor Juckamas Laohavanich, Ph.D}
{คณบดีคณะวิศวกรรมศาสตร์}
{Dean of The Faculty of Engineering}

\deanGrad
{ผศ. ดร. พลเดช เชาวรัตน์}
{ผู้ช่วยศาสตราจารย์ ดร.พลเดช เชาวรัตน์}
{Asst. Prof. Pondej Chaowarat, Ph.D}
{Assistant Professor Pondej Chaowarat, Ph.D}
{คณบดีบัณฑิตวิทยาลัย}
{Dean of Graduate School}

\degree
{ปรัชญาดุษฎีบัณฑิต}
{Doctor of Philosophy}
{วิศวกรรมไฟฟ้าและคอมพิวเตอร์}
{Electrical and Computer Engineering}
{คณะวิศวกรรมศาสตร์}
{The Faculty of Engineering}

\academicyear
{กรกฎาคม}
{July}
{2569}
{2026}

\abstract
{พิมพ์เนื้อหาบทคัดย่อภาษาไทยตรงนี้ความยาวประมาณ 300--500 คำ สรุปวัตถุประสงค์ วิธีการดำเนินการ ผลการทดลองและข้อสรุปของงานวิจัย}
{คำสำคัญที่ 1, คำสำคัญที่ 2, คำสำคัญที่ 3, คำสำคัญที่ 4}
{Type the English abstract content here, approximately 300--500 words. Summarize the research objectives, methodology, experimental results,and conclusions of the research.}
{Keyword 1, Keyword 2, Keyword 3, Keyword 4}

\acknowledgement
{กิตติกรรมประกาศกิตติกรรมประกาศกิตติกรรมประกาศกิตติกรรมประกาศกิตติกรรมประกาศกิตติกรรมประกาศกิตติกรรมประกาศกิตติกรรมประกาศ}
%
{acknowledgementacknowledgementacknowledgementacknowledgement acknowledgementacknowledgement acknowledgement}

%% ===================================================================
%%  ประวัติผู้เขียน (Biography)
%% ===================================================================
\biographyAuthor
{Chatmongkol Areeyat}                      %Name
{18th July 2002}                           %Date of birth   
{Nongbuarawae Hospital}                    %Place of birth
{                                          %Address
  292 Huanong, Nongbuarawae Sub-district,   
  Nongbuarawae District, Chaiyaphum, 36250
}
{Engineer}                                 %Position
{Asefa Public Company Limited}             %Place of work
{                                          %Education
  2024 – Bachelor of Engineering (B.Eng.) in Electrical Engineering, Mahasarakham University \newline
  2025 – Master of Engineering (M.Eng.) in Electrical and Computer Engineering, Mahasarakham University \newline
  2027 – Doctor of Philosophy (Ph.D.) in Electrical and Computer Engineering, Mahasarakham University 
}
{                                          %Research grants & awards
  - Areeyat, C., Audomsi, S., Obma, J., Yang, X., \& Sa-ngiamvibool, W. (2025). Load frequency control of multi-source power system using PID+DD controller based on chess algorithm. Bulletin of Electrical Engineering and Informatics, 14(6), 4900-4912. \newline
  - Areeyat, C., Audomsi, S., Obma, J., Yang, X., \& Sa-Ngiamvibool, W. (2025). Chess Optimizer for Load Frequency Control of Three-Area Multi-Source Renewable Energy Based on PID Plus Second Order Derivative Controller. International Journal of Robotics and Control Systems, 5(3), 2053-2067. \newline
  - Audomsi, S., Areeyat, C., Obma, J., \& Sa-ngiamvibool, W. (2026). Chess-Optimized Fuzzy-2DOF-PID Controller for Load Frequency Stabilization in Interconnected Multi-Area Renewable Grids. IEEE Access. \newline
  - Audomsi, S., Areeyat, C., Obma, J., Yang, X., \& Sa-ngiamvibool, W. (2025). An ALO-Tuned 2DOF-PID Controller for Enhanced Frequency Regulation in Hybrid PV–Thermal Power Systems. Engineering, Technology \& Applied Science Research, 15(5), 26517-26522. \newline
  - Pum, V., Audomsi, S., Areeyat, C., Luo, J., \& Chaiduangsri, N. (2026). Optimization-Based Tuning of a Cascaded PDN–PI Controller for Frequency Regulation in Solar PV–Powered Thermal Systems. Engineering Access, 12(1), 160-173. \newline
  - Riyawong, J., Audomsi, S., Phakditha, W., Areeyat, C.,\newline Pakdeesuwan, K., Ukumphan, S., \& Wattana, S. (2025). Hybrid PV-Reheat Thermal Power System Automatic Generation Control Using PIDD Controller Based on Hippopotamus Algorithm. Indochina Applied Sciences, 14(3), 263235-263235.
}
{}                                         %Research output